\section{Overview}
The first two chapters of this dissertation will focus on introducing and defining general concepts in \gls{rl} as a whole and in the delayed literature in particular.

\begin{itemize}
    \item \Cref{chap:prelim} introduces the main mathematical tools and algorithmic material used in the following chapters. 
    In particular, the \gls{mdp} model is introduced as well as the general framework of  \gls{rl}. 
    \item \Cref{chap:delay} will delve into the problem of delays as an extension of the original \gls{rl} framework described in the first chapter.
    Notations and terminology related to this topic will be introduced here.
    It will be followed by an extensive bestiary of the delays encountered in the literature or in practical applications.
    Finally, the literature on delays in \gls{rl} and related fields will be discussed. 
\end{itemize}

\noindent These introductory chapters are followed by four chapters where the original contribution of this dissertation is exposed. Each chapter will share the same structure, which will be as follows. A first section will introduce the problem formally and detail our proposed solution. The second section provides a theoretical analysis to better grasp the specificities of the problem and the properties of the solution.
Lastly, a third section will support the claims of the first two sections with a thorough empirical analysis. 

\begin{itemize}
    \item \Cref{chap:belief_based} focuses on the constant delay in state observation and action execution. A probabilistic approach called \emph{belief representation network} will be proposed in which the agent predicts its near future to account for the delay. 
    The theoretical analysis will provide a better understanding of the problem posed by this type of delay and its negative impact on performance.
    The content of this chapter can be found in \cite{liotet2021learning}.
    \item \Cref{chap:imitation_undelayed} places itself in the same framework as in the previous chapter. 
    A simple solution--{DIDA}--is proposed where the delayed agent learns to imitate an undelayed expert. 
    Despite its simplicity, we demonstrate great theoretical guarantees for DIDA in smooth environments. 
    The solution is extended to new frameworks, including non-integer and stochastic delays.
    The content of this chapter is inspired by \cite{liotet2022delayed}.
    \item \Cref{chap:lifelong} is again set in the same constant delay framework and proposes yet another approach to it.
    The solution considers a memoryless agent, i.e., blind to the delay.
    From the point of view of the agent, the process becomes non-stationary and we, therefore, propose to design a non-stationary policy to adapt to the evolving dynamics. 
    We do so by optimising an estimate of the agent's future performance and call this approach POLIS.
    This chapter provides theoretical ground for the design of POLIS, by studying in particular the bias and variance of the estimator of the future performance. 
    The content of this chapter is taken from \cite{liotet2022lifelong}.
    \item \Cref{chap:multi_action_delay} considers a delay framework that includes the constant delay one.
    In it, the execution of an action can be spread over multiple future steps, thus creating a multiple execution delay.
    The new setting's properties are thoroughly analysed, and a particular focus is put on understanding the effect of the delay's distribution on the best performance the agent could get. 
    The material in this chapter has not been published in any other location.
\end{itemize}

Finally, \Cref{chap:conclusion} concludes the dissertation by recalling the main achievements and takeaways, as well as limitations.
Importantly, for each result, the reader will be directed to the part of the dissertation where the result is presented. 
This chapter will also be an opportunity to discuss potential research directions.

Further proofs, experiment details and empirical results can be found in the appendix, in \Cref{app:proofs_results}.